\section{2 材料与方法}

\subsection{3.1 研究区域概况}

本研究选取某校园污水管网作为研究对象。该校园占地面积约X公顷，常住人口约X万人，日均污水排放量约X m³/d。校园功能区主要包括教学区、生活区、餐饮区和运动区，各功能区的污水通过支管汇入主管道后排出校园。

\subsection{3.2 采样方案}

根据管网布局和功能区分布，在主管道上设置了X个采样点，分别位于教学区(A1-A3)、生活区(B1-B3)、餐饮区(C1-C3)和管口出口(D)。采样时间涵盖冬季(2024年X月)和春季(2025年X月)两个季节，每次采样在各点同步采集固相、液相和气相样品。

\subsection{3.3 分析方法}

气相分析：采用便携式气体检测仪测定管道内CH₄、CO₂、O₂和VOCs浓度。检测前进行仪器校准，每个采样点重复测定3次取平均值。

液相分析：采集管道内污水水样，经0.45μm滤膜过滤后测定溶解性指标。总有机碳(TOC)采用TOC分析仪测定(HJ 501-2009)；化学需氧量(COD)采用重铬酸盐法测定(GB 11914-89)；总氮(TN)采用碱性过硫酸钾消解紫外分光光度法(HJ 636-2012)；铵态氮(NH₄⁺-N)采用纳氏试剂分光光度法(HJ 535-2009)。

固相分析：采集管道底部沉积物样品，自然风干后研磨过筛。固相总碳和有机碳采用元素分析仪测定。

\subsection{3.4 数据处理与统计分析}

采用Python 3.11进行数据处理和统计分析。描述性统计计算均值、标准差、变异系数等指标。正态性检验采用Shapiro-Wilk检验(p>0.05为正态)。组间差异分析：正态数据采用独立样本t检验，非正态数据采用Mann-Whitney U检验。相关性分析采用Pearson相关系数。降维分析采用主成分分析(PCA)，聚类分析采用层次聚类分析(HCA)。统计显著性水平设为p<0.05。

\subsection{3.5 碳平衡计算方法}

碳平衡基于质量守恒原理，计算公式为：

C\_input = C\_output + C\_accumulation + C\_loss

其中，C\_input为输入碳量，C\_output为输出碳量，C\_accumulation为碳储量变化，C\_loss为碳转化损失量。碳在固-液-气三相中的分配比例按各相碳含量占总碳含量的百分比计算。